\documentclass[11pt]{article}
\usepackage[utf8]{inputenc}
\usepackage{amsmath,amssymb}
\usepackage[margin=1in]{geometry}
\usepackage{hyperref}
\emergencystretch=3em

\title{The regular case: $S_{1/2}$'s Gaussian form and the half-integer ladder}
\author{Claude, with Daniel Murfet}
\date{2026-09-06}

\begin{document}
\maketitle
\sloppy

\begin{abstract}
Unit 10 of the grammar-paper \S4 autoformalisation, formalising Example \texttt{ex:fluctuation\_half}
--- the regular-model case $\lambda=1/2$. Two facts: (a) the \emph{half-integer ladder}
$\partial_a^n S_{1/2}(a)=\beta^n S_{(n+1)/2}(a)$, so every half-integer fluctuation function is an
$a$-derivative of the single function $S_{1/2}$; and (b) the \emph{Gaussian form}
$S_{1/2}(a)=\sqrt{2/\beta}\,e^{\beta a^2/4}\int_0^\infty e^{-(u-a\sqrt{\beta/2})^2/2}\,du$ from
completing the square. Zero \texttt{sorry}/\texttt{axiom}.
\end{abstract}

\section{Setting}
For regular statistical models the resolution is trivial and the RLCT is $\lambda=1/2$ on every chart,
so the whole fluctuation family collapses onto $S_{1/2}$. The paper records this in
\texttt{ex:fluctuation\_half}: $S_{1/2}$ has an error-function closed form, and all higher
$S_{(n+1)/2}$ are its $a$-derivatives. Both facts sit directly on the derivative ladder (unit~2) and
the kernel form (unit~7), needing no new machinery.

\section{The things formalised}
{\small
\begin{verbatim}
theorem fluctuation_half_iteratedDeriv (β : ℝ) (hβ : 0 < β) (n : ℕ) :
    iteratedDeriv n (fun a => fluctuation β (1 / 2) a)
      = fun a => β ^ n * fluctuation β (((n : ℝ) + 1) / 2) a

theorem fluctuation_half_gaussian (β a : ℝ) (hβ : 0 < β) :
    fluctuation β (1 / 2) a
      = Real.sqrt (2 / β) * Real.exp (β * a ^ 2 / 4)
        * ∫ u in Set.Ioi 0, Real.exp (-(u - a * Real.sqrt (β / 2)) ^ 2 / 2)
\end{verbatim}
}

\section{Proof strategy}
\textbf{Half-integer ladder.} Induction on $n$. \texttt{iteratedDeriv\_succ} peels one derivative;
the induction hypothesis turns the inner iterated derivative into $\beta^n S_{(n+1)/2}$, and
\texttt{deriv\_const\_mul} then \texttt{deriv\_fluctuation} (the ladder $S'_\lambda=\beta
S_{\lambda+1/2}$, valid since $(n+1)/2>0$) advance it to $\beta^{n+1}S_{(n+2)/2}$; index arithmetic
$(n+1)/2+1/2=(n+2)/2$ and $\beta^n\beta=\beta^{n+1}$ close it.

\textbf{Gaussian form.} Start from the kernel identity of unit~7 at $\lambda=1/2$:
$S_{1/2}(a)=2^{1/2}\beta^{-1/2}\int_0^\infty u^{0}\,e^{-u^2/2+a\sqrt{\beta/2}\,u}\,du$. The constant is
$\sqrt{2/\beta}$ (\texttt{Real.sqrt\_div'} plus $2^{1/2}=\sqrt2$, $\beta^{-1/2}=(\sqrt\beta)^{-1}$).
Pull the constant $e^{\beta a^2/4}$ out of the integral (\texttt{integral\_const\_mul} backwards) and
match integrands pointwise: $u^{0}=1$, and completing the square
$-u^2/2+a\sqrt{\beta/2}\,u=\beta a^2/4-(u-a\sqrt{\beta/2})^2/2$ using $(\sqrt{\beta/2})^2=\beta/2$.

\section{Roadblocks and resolutions}
Almost none:
\begin{itemize}
\item The completed-square identity is an equality with a $\sqrt{}$ that must be squared away.
\texttt{nlinarith} does not prove equalities; \texttt{linear\_combination (a\^{}2/2) * hs} (with
\texttt{hs :} $\sqrt{\beta/2}^2=\beta/2$) is the right tool --- the difference of the two sides is
exactly $\tfrac{a^2}{2}(\sqrt{\beta/2}^2-\beta/2)$.
\item \texttt{setIntegral\_congr\_fun} hands an applied-lambda goal; \texttt{beta\_reduce} before the
rewrites.
\item Mathlib has no \texttt{erf}, so the Gaussian \emph{tail} $\int_0^\infty e^{-(u-c)^2/2}\,du$ is
left un-evaluated; the statement is faithful to \texttt{ex:fluctuation\_half} up to that final
notation.
\end{itemize}

\section{What was Mathlib, what was new}
Mathlib supplied \texttt{iteratedDeriv\_succ}, \texttt{deriv\_const\_mul}, \texttt{integral\_const\_mul},
\texttt{Real.sqrt\_div'}/\texttt{sq\_sqrt}, and \texttt{linear\_combination}. New is the pairing of
unit~2's ladder and unit~7's kernel form into the two regular-case statements --- in particular the
observation that the $\lambda=1/2$ kernel integrand, after $u^{0}=1$ and one completed square, is a
shifted Gaussian.

\section{Lessons}
\texttt{linear\_combination} is the clean way to discharge a polynomial equality that holds only after a
square-root identity: name $\sqrt{c}^2=c$ and give the coefficient, rather than reaching for
\texttt{nlinarith} (which targets inequalities). And a specialised value of a general closed form (here
$S_{1/2}$ from the kernel form) is often cheaper to reshape directly than to re-derive from the integral
definition.

\section{Follow-ups}
The remaining \S4 items are operator-algebraic --- the ladder operators $b,b^\dagger$ with
$[b,b^\dagger]=1$, the number operator, and the log-insertion generating function --- which want a
Weyl-algebra framework, plus the order-derivative identity
$\partial_\lambda^k S_\lambda=\int(\log t)^k t^{\lambda-1}e^{\cdots}$ (differentiation under the
integral in the order parameter) that underlies \texttt{lem:log\_shift}/\texttt{prop:LogODE}.
\end{document}
